% Tag-Set Matching Ontology System
% A Multi-Layer Semantic Graph for Self-Organization of Heterogeneous Repositories
% Skeleton placeholder: full manuscript is maintained separately and will be merged here.

\documentclass[11pt,a4paper]{article}
\usepackage{amsmath}
\usepackage{amssymb}
\usepackage{hyperref}
\usepackage{booktabs}
\usepackage{geometry}
\geometry{margin=25mm}

\title{Tag-Set Matching Ontology System\\
A Multi-Layer Semantic Graph for Self-Organization of Heterogeneous Repositories}
\author{BONSAI Research}
\date{2026}

\begin{document}
\maketitle

\begin{abstract}
This paper proposes a Tag-Set Matching Ontology System for connecting heterogeneous repositories, domains, capabilities, agents, and workflows through tag sets treated as semantic signatures. The system models semantic similarity, containment, and transformation as a multi-layer graph and targets self-organization of heterogeneous software and knowledge assets.
\end{abstract}

\section{Introduction}
% TODO: insert full manuscript.
\section{Basic Hypothesis}
% TODO
\section{Multi-Layer Model}
% TODO
\section{Tag-Set Matching}
% TODO
\section{Multi-Layer Tags}
% TODO
\section{Cross-Domain Cluster: Dots}
% TODO
\section{Semantic Transformation}
% TODO
\section{Declaration and Inference}
% TODO
\section{Repository Graph}
% TODO
\section{repo2agent}
% TODO
\section{Self-Organization}
% TODO
\section{Implementation Architecture}
% TODO
\section{Minimal Data Model}
% TODO
\section{Discussion}
% TODO
\section{Expected Effects}
% TODO
\section{Conclusion}
% TODO

\appendix
\section{Core Concepts}
% TODO
\section{Application to BONSAI}
% TODO

\end{document}
